\chapter{Markov Chains}

\begin{center}\itshape
After five riffle shuffles a deck of $52$ cards is still far from random; after seven it is close,\\
and after ten nobody could tell. \upshape(after Bayer and Diaconis, 1992)
\end{center}

In 1913 Andrei Markov, to show that the law of large numbers does not need independence, took the first $20\,000$ letters of Pushkin's verse novel
\emph{Eugene Onegin}, wrote down the sequence of vowels and consonants, and counted. A vowel was followed by a vowel only $13\%$ of the time, a
consonant by a vowel $66\%$ of the time: the letters were far from independent, yet their frequencies settled down exactly as his theory predicted. It
was the first application of what are now called Markov chains, and they have since run on every computer, from the samplers of physicists to the
ranking of web pages.

A Markov chain is a random process that forgets its past: where it goes next depends only on where it is now. The
random walks of Chapter~7 are Markov chains, and so are a shuffled deck of cards, a web surfer clicking links, the
molecules in Ehrenfest's two urns, and the samplers that statisticians and physicists run on computers to explore
huge spaces.

The central theorem says that a finite chain that can go everywhere, and is not trapped in a cycle, forgets where it
started: its law converges to a unique stationary distribution. We prove it twice, by a coupling argument of
Doeblin that gives an explicit rate, and through the ergodic theorem, which says that time averages along one path
converge to averages over the stationary law. Then we ask the quantitative question that matters in practice:
\emph{how many steps} until the chain is close to stationary? For the walk on a cycle, Fourier analysis gives the answer,
of order $n^2$. For shuffling cards by random transpositions, a clever stopping rule gives $2n\ln n$ and a count of
untouched cards gives the lower bound $\frac12n\ln n$. For riffle shuffles, an exact formula of Bayer and Diaconis computes the
distance to a perfectly mixed deck after $k$ shuffles, and shows the sudden drop between $5$ and $8$ shuffles.

\begin{think}
A knight starts in a corner of a chessboard and makes random legal moves, each move chosen uniformly among those
available. How many moves does it take, on average, to come back to its corner? Section~\ref{sec:reversible} answers in two
lines, without solving any equation.
\end{think}

\section{Transition matrices}

Let $E$ be a finite or countable set, the \emph{state space}. A \emph{transition matrix} on $E$ is a matrix $P=(P(x,y))_{x,y\in E}$ with non-negative
entries and rows summing to $1$; $P(x,y)$ is the probability of jumping from $x$ to $y$.

\begin{definition}\label{def:markov}
A sequence $(X_n)_{n\ge0}$ of random variables with values in $E$ is a \emph{Markov chain} with initial law $\lambda$ and transition matrix
$P$ if, for all $n$ and all $x_0,\dots,x_n\in E$, the probability $P(X_0=x_0,\dots,X_n=x_n)$ is
\[
\lambda(x_0)P(x_0,x_1)\cdots P(x_{n-1},x_n).
\]
We write $P_x$ and $\mathbb E_x$ when $\lambda=\delta_x$, that is, when the chain starts at $x$.
\end{definition}

Such a chain exists: list $E=\{e_1,e_2,\ldots\}$, take $X_0$ with law $\lambda$ and independent uniform variables $U_1,U_2,\ldots$ (admitted fact (A1)
of Chapter~5), and let $X_{n+1}=e_k$ when $\sum_{j<k}P(X_n,e_j)\le U_{n+1}<\sum_{j\le k}P(X_n,e_j)$. Then $P(X_{n+1}=y\mid X_0,\dots,X_n)=P(X_n,y)$, and the product
formula follows by induction. This is also how a computer simulates a chain.

Summing the product formula over $x_0,\dots,x_{n-1}$ gives $P(X_n=y)=(\lambda P^n)(y)$, where $\lambda$ is a row vector and $P^n$ the $n$-th matrix power:
the law at time $n$ is $\lambda P^n$, and $P^n(x,y)=P_x(X_n=y)$. The relation $P^{m+n}=P^mP^n$ is the \emph{Chapman--Kolmogorov equation}.

\begin{proposition}[Markov property]\label{prop:markovprop}
Let $(X_n)$ be a Markov chain with transition matrix $P$, $m\ge0$ and $x\in E$ with $P(X_m=x)>0$. Conditionally on $X_m=x$, the
process $(X_{m+n})_{n\ge0}$ is a Markov chain started at $x$ with transition matrix $P$, and it is independent of $(X_0,\dots,X_m)$.
\end{proposition}
\begin{proof}
For any $x_0,\dots,x_{m-1}$ and $y_1,\dots,y_n$, let $B=\{X_0=x_0,\dots,X_{m-1}=x_{m-1},X_m=x\}$ and $C=\{X_{m+1}=y_1,\dots,X_{m+n}=y_n\}$. The product formula gives
\[
\begin{aligned}P(B\cap C)&=P(B)\eqbr\cdot P(x,y_1)P(y_1,y_2)\cdots P(y_{n-1},y_n).\end{aligned}
\]
Dividing by $P(X_m=x)$ and summing over $x_0,\dots,x_{m-1}$ in a set $A$ of paths shows that the conditional probability of
$\{(X_0,\dots,X_m)\in A\}\cap\{X_{m+1}=y_1,\dots,X_{m+n}=y_n\}$ is the product of the conditional probability of the first event and of
$P_x(X_1=y_1,\dots,X_n=y_n)$.
\end{proof}

In practice we use it at random times $T$ by splitting according to the value of $T$, exactly as in Chapter~7: if an
event $B$ is determined by $X_0,\dots,X_m$ and contained in $\{X_m=x\}$, then the chain after time $m$ is, on $B$, a fresh chain from $x$,
independent of $B$.

\begin{example}\label{ex:examples}
(a) The simple random walk on a graph: $P(x,y)=\frac1{\deg x}$ if $y$ is a neighbour of $x$. (b) \emph{Ehrenfest's urns}: $N$ balls are
split between two urns, and at each step a uniformly chosen ball changes urn; the state is the number $x$ of balls in the
first urn, with $P(x,x-1)=\frac xN$ and $P(x,x+1)=\frac{N-x}N$. (c) \emph{Shuffling}: $E=\mathfrak S_n$, the orderings of a deck, and each step applies
a random permutation drawn from a fixed law $Q$: $P(\sigma,\tau\sigma)=Q(\tau)$. This is a \emph{random walk on the group} $\mathfrak S_n$.
\end{example}

\begin{figure}[ht]
\centering
\fig{ch08-chain}
\caption{A Markov chain on three states, drawn as a graph with the transition probabilities on the arrows. Its
stationary law is computed in the text.}
\label{fig:chain}
\end{figure}

\section{Irreducibility and periodicity}

A state $y$ is \emph{accessible} from $x$ if $P^n(x,y)>0$ for some $n\ge0$; the chain is \emph{irreducible} if every state is accessible from
every other one. The \emph{period} of $x$ is the greatest common divisor of $S_x=\{n\ge1:P^n(x,x)>0\}$, and $x$ is \emph{aperiodic} if its period is $1$. For
instance, the simple random walk on a bipartite graph, such as $\Z$ or a cycle of even length, has period $2$; a chain with
$P(x,x)>0$ at some state $x$ has period $1$ at $x$.

\begin{lemma}\label{lem:period}
In an irreducible chain all states have the same period.
\end{lemma}
\begin{proof}
Take $x,y$ and $a,b$ with $P^a(x,y)>0$ and $P^b(y,x)>0$. Then $P^{a+b}(x,x)\ge P^a(x,y)P^b(y,x)>0$, and for $n\in S_y$, $P^{a+n+b}(x,x)\ge P^a(x,y)P^n(y,y)P^b(y,x)>0$. So the
period $d_x$ of $x$ divides $a+b$ and $a+b+n$, hence every $n\in S_y$, hence $d_y$. By symmetry $d_y\mid d_x$.
\end{proof}

\begin{lemma}\label{lem:semigroup}
Let $S\subseteq\{1,2,\ldots\}$ be closed under addition, with greatest common divisor $1$. Then $S$ contains every integer from some
point on.
\end{lemma}
\begin{proof}
The gcd of $S$ is already the gcd of finitely many elements $s_1,\dots,s_m$ (the gcd of the first $k$ elements decreases and
stabilises). By Bézout, $1=\sum_ic_is_i$ with $c_i\in\Z$. Grouping the positive and the negative terms, $1=a-b$ where $a$ and $b$ are sums of
elements of $S$, so $a\in S$ and $b\in S\cup\{0\}$. If $b=0$, then $1\in S$ and $S=\{1,2,\ldots\}$. Otherwise let $n\ge b^2$ and write $n=qb+\rho$ with $0\le\rho<b\le q$.
Then $n=qb+\rho(a-b)=(q-\rho)b+\rho a$, a sum of elements of $S$ (with $q-\rho\ge1$), so $n\in S$.
\end{proof}

\begin{proposition}\label{prop:primitive}
If $P$ is irreducible and aperiodic on a finite set $E$, there is $N$ such that $P^n(x,y)>0$ for all $x,y\in E$ and all $n\ge N$.
\end{proposition}
\begin{proof}
The set $S_y$ is closed under addition, since $P^{m+n}(y,y)$ is at least $P^m(y,y)P^n(y,y)$, and it has gcd $1$; by Lemma~\ref{lem:semigroup} it contains all $n\ge N_y$. For $x,y$ pick $m=m_{xy}$ with
$P^m(x,y)>0$; then for $n\ge m+N_y$,
\[
P^n(x,y)\ge P^m(x,y)P^{n-m}(y,y)>0.
\] Take $N=\max_{x,y}(m_{xy}+N_y)$, a finite maximum.
\end{proof}

For a finite irreducible chain, let $T_y=\min\{n\ge1:X_n=y\}$ be the time of the first visit to $y$ after time $0$ (the first \emph{return} if
the chain starts at $y$).

\begin{lemma}\label{lem:hitting}
If $P$ is irreducible on a finite set $E$, there are $K\ge1$ and $\alpha>0$ such that $P_x(T_y>kK)\le(1-\alpha)^k$ for all $x,y$ and $k$. In particular
$\mathbb E_x[T_y^p]<\infty$ for every $p\ge1$.
\end{lemma}
\begin{proof}
For each pair $x,y$ choose a path of positive probability from $x$ to $y$ with at least one step and at most $K=|E|$ steps (remove the
loops of any path; for $x=y$, go to a neighbour and come back). Let $\alpha>0$ be the smallest of the probabilities of these
paths. Whatever the state at time $jK$, the chain then follows the chosen path to $y$ during the next $K$ steps with probability
at least $\alpha$, by the Markov property; so $P_x(T_y>(j+1)K\mid T_y>jK)\le1-\alpha$, and the bound follows by induction. Then
$\mathbb E_x[T_y^p]\le\sum_kK^p(k+1)^pP_x(T_y>kK)<\infty$.
\end{proof}

\section{Stationary distributions}

A probability $\pi$ on $E$ is \emph{stationary} for $P$ if $\pi P=\pi$: a chain started from $\pi$ has law $\pi$ at every time.

\begin{theorem}[existence and uniqueness]\label{thm:stationary}
An irreducible chain on a finite set $E$ has exactly one stationary distribution $\pi$. It charges every state, and
\[
\pi(x)=\frac1{\mathbb E_x[T_x]}\qquad\text{(Kac's formula).}
\]
\end{theorem}
\begin{proof}
\emph{Existence.} Fix $x$ and let $\mu(y)=\mathbb E_x\bigl[\sum_{n=0}^{T_x-1}\mathbf 1_{X_n=y}\bigr]=\sum_{n\ge0}P_x(X_n=y,\ n<T_x)$, the expected number of visits to $y$ during an
excursion from $x$. Then $\mu(x)=1$ and $\sum_y\mu(y)=\mathbb E_x[T_x]<\infty$ by Lemma~\ref{lem:hitting}. The event $\{X_n=y,\ n<T_x\}$ is determined by $X_0,\dots,X_n$,
so by the Markov property $P_x(X_n=y,\ n<T_x,\ X_{n+1}=z)=P_x(X_n=y,\ n<T_x)P(y,z)$. Summing over $y$ and $n$,
\[
\begin{aligned}(\mu P)(z)&=\sum_{n\ge0}P_x(X_{n+1}=z,\ n<T_x)\eqbr=\sum_{m\ge1}P_x(X_m=z,\ m\le T_x).\end{aligned}
\]
For $z\ne x$ the terms with $m=T_x$ vanish and the term $m=0$ of $\mu(z)$ is $0$, so $(\mu P)(z)=\mu(z)$. For $z=x$, $X_m=x$ with $m\le T_x$ happens only for $m=T_x$,
so $(\mu P)(x)=P_x(T_x<\infty)=1=\mu(x)$. So $\pi_x=\mu/\mathbb E_x[T_x]$ is stationary, and $\pi_x(x)=\frac1{\mathbb E_x[T_x]}$.

\emph{Uniqueness.} If $Ph=h$ for a function $h$ on $E$ (a column vector), let $x$ be a point where $h$ is maximal. Since $h(x)=\sum_yP(x,y)h(y)$
is an average of values $\le h(x)$, $h(y)=h(x)$ whenever $P(x,y)>0$; by induction $h=h(x)$ on every state accessible from $x$, that is,
everywhere. So the kernel of $P-I$ is the line of constants, $P-I$ has rank $|E|-1$, and so has its transpose: the solutions of
$\pi P=\pi$ form a line, which contains exactly one probability. Since $\pi_x$ is stationary for every choice of $x$, all the $\pi_x$ are
equal to it, which gives Kac's formula for every state. Finally $\pi(x)=\frac1{\mathbb E_x[T_x]}>0$.
\end{proof}

For the chain of Figure~\ref{fig:chain}, solving $\pi P=\pi$ gives $\pi=(\frac4{11},\frac4{11},\frac3{11})$, so a chain started at $1$ returns to it after
$\frac{11}4$ steps on average.

\section{Convergence and coupling}

To say how close two laws are we use the \emph{total variation distance}
\[
\begin{aligned}\|\mu-\nu\|&=\max_{A\subseteq E}|\mu(A)-\nu(A)|\eqbr=\frac12\sum_{x\in E}|\mu(x)-\nu(x)|.\end{aligned}
\]
(The two expressions agree: the maximum is attained at $A=\{\mu>\nu\}$, and $\mu(A)-\nu(A)=\nu(A^c)-\mu(A^c)$ since both laws have mass $1$.) It is the
largest error one can make on the probability of an event by using $\nu$ instead of $\mu$.

\begin{lemma}[coupling inequality]\label{lem:coupling}
If $X$ has law $\mu$ and $Y$ has law $\nu$, defined on the same probability space, then $\|\mu-\nu\|\le P(X\ne Y)$.
\end{lemma}
\begin{proof}
$\mu(A)-\nu(A)=P(X\in A)-P(Y\in A)\le P(X\in A,\ Y\notin A)\le P(X\ne Y)$, and likewise with $\mu$ and $\nu$ exchanged.
\end{proof}

A pair $(X,Y)$ with given marginals is a \emph{coupling}. Building couplings in which $X=Y$ with high probability is the main way to
show that two laws are close, and the next proof is the first instance: Doeblin's idea (1938) is to let two copies of
the chain meet.

\begin{lemma}\label{lem:contraction}
For any transition matrix $Q$ and probabilities $\mu,\nu$: $\|\mu Q-\nu Q\|\le\|\mu-\nu\|$. If moreover $Q(x,z)\ge\delta$ for all $x,z\in E$, where $E$ is
finite with $m$ elements, then $\|\mu Q-\nu Q\|\le(1-m\delta)\|\mu-\nu\|$.
\end{lemma}
\begin{proof}
First,
\[
\begin{aligned}{}&\sum_z\Bigl|\sum_x(\mu(x)-\nu(x))Q(x,z)\Bigr|\eqbr\le\sum_x|\mu(x)-\nu(x)|\sum_zQ(x,z)\eqbr=\sum_x|\mu(x)-\nu(x)|.\end{aligned}
\]
For the second part, $\theta=m\delta\le1$; if $\theta=1$, $Q$ has all
rows uniform and $\mu Q=\nu Q$. Otherwise $Q=\theta U+(1-\theta)R$, where $U$ has all rows equal to the uniform law and $R=\frac{Q-\theta U}{1-\theta}$ is a
transition matrix (non-negative entries since $Q(x,z)\ge\delta=\frac\theta m$). As $\mu U=\nu U$, $\mu Q-\nu Q=(1-\theta)(\mu R-\nu R)$, and the first part concludes.
\end{proof}

The decomposition has a coupling meaning: at each step, with probability $\theta$, both copies of the chain jump to the same
uniform state, after which they can move together forever.

\begin{theorem}[convergence theorem]\label{thm:convergence}
Let $P$ be irreducible and aperiodic on a finite set $E$, with stationary law $\pi$. There are $C$ and $\rho<1$ such that, for every initial
law $\lambda$ and every $n$,
\[
\|\lambda P^n-\pi\|\le C\rho^n.
\]
\end{theorem}
\begin{proof}
By Proposition~\ref{prop:primitive}, $Q=P^N$ has all entries $\ge\delta$ for some $N$ and $\delta>0$; let $\theta=|E|\delta$. Since $\pi Q=\pi$, Lemma~\ref{lem:contraction}
gives $\|\lambda P^{kN}-\pi\|=\|\lambda Q^k-\pi Q^k\|\le(1-\theta)^k$, as $\|\lambda-\pi\|\le1$. For $n=kN+j$ with $0\le j<N$, the first part of the lemma gives
$\|\lambda P^n-\pi\|=\|(\lambda P^j)P^{kN}-\pi\|\le(1-\theta)^k$. Take $\rho=(1-\theta)^{1/N}$ and $C=(1-\theta)^{-1}$.
\end{proof}

Without aperiodicity there is no convergence: the walk on a cycle of even length alternates between even and odd
points. The cure is to make the chain \emph{lazy}, replacing $P$ by $\frac12(I+P)$: it stays put with probability $\frac12$, has the same
stationary law, and is aperiodic (Exercise~\ref{ex:lazy}).

To measure the speed of convergence, let
\[
\begin{gathered}d(t)=\max_{x\in E}\|P^t(x,\cdot)-\pi\|,\qq t_{\mathrm{mix}}(\varepsilon)=\min\{t:d(t)\le\varepsilon\},\end{gathered}
\]
and $t_{\mathrm{mix}}=t_{\mathrm{mix}}(\frac14)$, the \emph{mixing time}. As every law is a mixture of point masses, $\|\lambda P^t-\pi\|\le d(t)$ for every $\lambda$ (convexity of the
norm).

\begin{lemma}\label{lem:submult}
Let $\bar d(t)=\max_{x,y}\|P^t(x,\cdot)-P^t(y,\cdot)\|$. Then $d(t)\le\bar d(t)\le2d(t)$ and $\bar d(s+t)\le\bar d(s)\bar d(t)$. Consequently $d(kt_{\mathrm{mix}})\le2^{-k}$.
\end{lemma}
\begin{proof}
$\bar d\le2d$ by the triangle inequality. Since $\pi=\sum_y\pi(y)P^t(y,\cdot)$, $\|P^t(x,\cdot)-\pi\|\le\sum_y\pi(y)\|P^t(x,\cdot)-P^t(y,\cdot)\|\le\bar d(t)$. For submultiplicativity,
let $\mu=P^s(x,\cdot)$, $\nu=P^s(y,\cdot)$ and $\varepsilon=\|\mu-\nu\|$. If $\varepsilon>0$, the positive and negative parts of $\mu-\nu$ both have mass $\varepsilon$, so $\mu-\nu=\varepsilon(\alpha-\beta)$
with $\alpha,\beta$ probabilities. With $\Delta_{uv}=P^t(u,\cdot)-P^t(v,\cdot)$,
\[
\|(\mu-\nu)P^t\|=\varepsilon\Bigl\|\sum_{u,v}\alpha(u)\beta(v)\Delta_{uv}\Bigr\|\le\varepsilon\,\bar d(t),
\]
and $\varepsilon\le\bar d(s)$. Finally $d(kt)\le\bar d(kt)\le\bar d(t)^k\le(2d(t))^k$, and $2d(t_{\mathrm{mix}})\le\frac12$.
\end{proof}

So once the distance is below $\frac14$, it halves at least every $t_{\mathrm{mix}}$ further steps: the mixing time is the one number that
matters.

\section{The ergodic theorem}

The convergence theorem is about the law of $X_n$. The ergodic theorem is about one trajectory: the fraction of time it
spends in a state converges to the stationary probability of that state. It holds without aperiodicity.

\begin{theorem}[ergodic theorem]\label{thm:ergodic}
Let $(X_n)$ be an irreducible Markov chain on a finite set $E$, with any initial law, and let $\pi$ be its stationary law. For
every function $f:E\to\R$, almost surely,
\[
\frac1n\sum_{k=0}^{n-1}f(X_k)\xrightarrow[n\to\infty]{}\sum_{x\in E}\pi(x)f(x).
\]
\end{theorem}
\begin{proof}
Fix a state $x$. Let $R_0$ be the first time $\ge0$ at which $X_{R_0}=x$, and $R_{j+1}=\min\{n>R_j:X_n=x\}$; all are finite almost surely, by
Lemma~\ref{lem:hitting} applied repeatedly. The $j$-th \emph{excursion} is $W_j=(X_{R_j},\dots,X_{R_{j+1}-1})$, a finite word starting with $x$.

\emph{The excursions are independent and identically distributed.} For words $w_0,\dots,w_j$ that start with $x$ and do not
contain $x$ elsewhere,
the probability that $W_0=w_0,\dots,W_j=w_j$ is
\[
P(W_0=w_0,\dots,W_{j-1}=w_{j-1})\cdot P_x(W_0=w_j).
\]
Indeed, split according to the value $m$ of $R_j$: the event $\{R_j=m,\ W_0=w_0,\dots,W_{j-1}=w_{j-1}\}$ is determined by $X_0,\dots,X_m$ and
contained in $\{X_m=x\}$, so by Proposition~\ref{prop:markovprop} the chain after time $m$ is a fresh chain from $x$, independent of it, and $W_j=w_j$
is the event that this fresh chain's first excursion is $w_j$. By induction the probability is $\prod_iP_x(W_0=w_i)$.

\emph{Strong law.} Let $\tau_j=R_{j+1}-R_j$ and $Z_j=\sum_{k=R_j}^{R_{j+1}-1}f(X_k)$, functions of $W_j$; the pairs $(\tau_j,Z_j)$ are i.i.d., with $|Z_j|\le\|f\|\tau_j$ where
$\|f\|=\max|f|$, and $\tau_j$ has the law of $T_x$ under $P_x$, so all their moments are finite (Lemma~\ref{lem:hitting}). By the strong law
(Theorem~\ref{thm:slln4}), almost surely
\[
\begin{gathered}\frac{R_j-R_0}j=\frac1j\sum_{i<j}\tau_i\to\mathbb E_x[T_x],\qq\frac1j\sum_{i<j}Z_i\to\mathbb E[Z_0]=\sum_y\mu(y)f(y),\end{gathered}
\]
with $\mu$ as in the proof of Theorem~\ref{thm:stationary}, so $\mathbb E[Z_0]=\mathbb E_x[T_x]\sum_y\pi(y)f(y)$.

\emph{Sandwich.} For $n>R_0$ let $J_n$ be the number of $j\ge1$ with $R_j\le n$, so $R_{J_n}\le n<R_{J_n+1}$, and $J_n\to\infty$. Then
\[
\begin{aligned}\Bigl|\sum_{k<n}f(X_k)-\sum_{i<J_n}Z_i\Bigr|&\le\|f\|\bigl(R_0+\tau_{J_n}\bigr).\end{aligned}
\]
Since $\frac{R_{J_n}}{J_n}$ and $\frac{R_{J_n+1}}{J_n}$ both tend to $\mathbb E_x[T_x]$, $\frac n{J_n}\to\mathbb E_x[T_x]$ and $\frac{\tau_{J_n}}n\le\frac{R_{J_n+1}-R_{J_n}}{J_n}\cdot\frac{J_n}n\to0$. Dividing by $n$,
\[
\begin{aligned}\frac1n\sum_{k<n}f(X_k)&=\frac{J_n}n\cdot\frac1{J_n}\sum_{i<J_n}Z_i+o(1)\eqbr\to\frac{\mathbb E[Z_0]}{\mathbb E_x[T_x]}=\sum_y\pi(y)f(y).\end{aligned}
\]
\end{proof}

With $f=\mathbf 1_{\{y\}}$: the fraction of time spent at $y$ tends to $\pi(y)$, almost surely. This is what makes Markov chain Monte Carlo work: to
compute an average over a huge space, run a chain with the right stationary law and average along its path.

\section{Reversibility and random walks on graphs}\label{sec:reversible}

\begin{proposition}[detailed balance]\label{prop:detailed}
If a probability $\pi$ satisfies $\pi(x)P(x,y)=\pi(y)P(y,x)$ for all $x,y$, then $\pi$ is stationary. The chain is then called \emph{reversible}:
started from $\pi$, the path $(X_0,\dots,X_n)$ has the same law as $(X_n,\dots,X_0)$.
\end{proposition}
\begin{proof}
Summing detailed balance over $x$,
\[
\begin{aligned}(\pi P)(y)&=\sum_x\pi(x)P(x,y)\eqbr=\sum_x\pi(y)P(y,x)=\pi(y).\end{aligned}
\]
Reversing a path, applying detailed balance $n$ times gives
\[
\begin{gathered}\pi(x_0)P(x_0,x_1)\cdots P(x_{n-1},x_n)\qq=\pi(x_n)P(x_n,x_{n-1})\cdots P(x_1,x_0).\end{gathered}
\]
\end{proof}

\begin{corollary}\label{cor:graph}
For the simple random walk on a finite connected graph with edge set $\mathcal E$, $\pi(x)=\frac{\deg x}{2|\mathcal E|}$, and the expected return time to $x$ is $\frac{2|\mathcal E|}{\deg x}$.
\end{corollary}
\begin{proof}
$\deg(x)P(x,y)=1=\deg(y)P(y,x)$ for neighbours $x,y$, and both sides are $0$ otherwise; $\sum_x\deg x=2|\mathcal E|$. Kac's formula gives the return time.
\end{proof}

\begin{figure}[ht]
\centering
\fig{ch08-knight}
\caption{The number of legal knight moves from each square of a chessboard; from the corner (orange) only the two blue squares can be reached. The degrees add up to $336$, so a random
knight returns to its corner after $\frac{336}2=168$ moves on average (Exercise~\ref{ex:knight}).}
\label{fig:knight}
\end{figure}

For reversible chains, the speed of convergence is governed by the eigenvalues of $P$, which are real. This is the link
between Markov chains and the linear algebra of Tome~2 (Jordan forms describe $P^n$ for any chain; for reversible
chains $P$ is even diagonalisable).

\begin{proposition}[spectral bound]\label{prop:spectral}
Let $P$ be irreducible and reversible with respect to $\pi$ on a finite set. Then $P$ has real eigenvalues $1=\lambda_1\ge\lambda_2\ge\dots\ge\lambda_m\ge-1$,
and with $\lambda_*=\max(|\lambda_2|,|\lambda_m|)$ and $\pi_{\min}=\min_x\pi(x)$,
\[
d(t)\le\frac{\lambda_*^t}{2\pi_{\min}}.
\]
\end{proposition}
\begin{proof}
Use the inner product $\langle f,g\rangle_\pi=\sum_x\pi(x)f(x)g(x)$ on functions. Detailed balance says exactly that $P$, acting on column vectors, is
self-adjoint: $\langle Pf,g\rangle_\pi=\sum_{x,y}\pi(x)P(x,y)f(y)g(x)=\langle f,Pg\rangle_\pi$. By the spectral theorem (Tome~0, Theorem~7.7, in this inner product)
there is an orthonormal basis $f_1,\dots,f_m$ of eigenvectors, with real eigenvalues; $f_1=1$ with $\lambda_1=1$, and $|\lambda_j|\le1$ since $P$ is an
average operator ($\max|Pf|\le\max|f|$). The function $\frac{\mathbf 1_y}{\pi(y)}$ has coordinates $\langle\frac{\mathbf 1_y}{\pi(y)},f_j\rangle_\pi=f_j(y)$, so $\frac{\mathbf 1_y}{\pi(y)}=\sum_jf_j(y)f_j$ and
\[
\begin{aligned}\frac{P^t(x,y)}{\pi(y)}&=P^t\Bigl(\frac{\mathbf 1_y}{\pi(y)}\Bigr)(x)=\sum_j\lambda_j^tf_j(x)f_j(y)\eqbr=1+\sum_{j\ge2}\lambda_j^tf_j(x)f_j(y).\end{aligned}
\]
By Cauchy--Schwarz and $\sum_jf_j(x)^2=\frac1{\pi(x)}$ (the same expansion, evaluated at $y=x$), $\bigl|\frac{P^t(x,y)}{\pi(y)}-1\bigr|\le\frac{\lambda_*^t}{\sqrt{\pi(x)\pi(y)}}\le\frac{\lambda_*^t}{\pi_{\min}}$.
Multiply by $\frac{\pi(y)}2$ and sum over $y$.
\end{proof}

The \emph{spectral gap} $1-\lambda_*$ is thus the rate of convergence, and $t_{\mathrm{mix}}$ is at most of order $\frac{1}{1-\lambda_*}\ln\frac1{\pi_{\min}}$. For a
walk on a group, the eigenvalues are computed by Fourier analysis, and the next section does this on the cycle.

\section{The walk on a cycle}

Let $n\ge2$ and consider the \emph{lazy walk on $\Z/n\Z$}: from $x$ it stays at $x$ with probability $\frac12$ and moves to $x\pm1$ with probability $\frac14$
each. Its stationary law is uniform, by Corollary~\ref{cor:graph} or by symmetry. Tome~2 (Chapter~11) diagonalised such walks with the
characters $\chi_m(x)=e^{2i\pi mx/n}$ of $\Z/n\Z$; here is what this gives.

\begin{theorem}\label{thm:cycle}
For the lazy walk on $\Z/n\Z$ with $n\ge4$,
\[
\begin{gathered}\frac12e^{-2\pi^2t/n^2}\le d(t)\quad\text{for all }t,\qq d(t)\le e^{-\pi^2t/n^2}\quad\text{for }t\ge\frac{n^2}4.\end{gathered}
\]
In particular $\frac{\ln2}{2\pi^2}n^2\le t_{\mathrm{mix}}\le\bigl\lceil\frac{n^2}4\bigr\rceil$.
\end{theorem}
\begin{proof}
\emph{Fourier coefficients.} Let $\mu_t=P^t(0,\cdot)$; by symmetry $d(t)=\|\mu_t-U\|$, $U$ the uniform law. For a function $g$ on $\Z/n\Z$ let
$\hat g(m)=\sum_xg(x)e^{-2i\pi mx/n}$. One step multiplies the coefficient by $c_m=\frac12+\frac14(e^{2i\pi m/n}+e^{-2i\pi m/n})=\cos^2\frac{\pi m}n$, because
\[
\begin{aligned}\widehat{\mu_{t+1}}(m)&=\sum_x\Bigl(\tfrac12\mu_t(x)+\tfrac14\mu_t(x-1)\eqbr\qquad+\tfrac14\mu_t(x+1)\Bigr)e^{-2i\pi mx/n}\\&=c_m\,\widehat{\mu_t}(m).\end{aligned}
\]
So $\widehat{\mu_t}(m)=c_m^t$, while $\hat U(m)=0$ for $m\ne0$ and $\widehat{\mu_t}(0)=\hat U(0)=1$. We also need Parseval's identity
$\sum_x|g(x)|^2=\frac1n\sum_m|\hat g(m)|^2$, which follows by expanding the right side and using $\sum_me^{2i\pi m(x-y)/n}=n\mathbf 1_{x=y}$ (a geometric sum).

\emph{Upper bound.} By Cauchy--Schwarz and Parseval,
\[
\begin{aligned}\|\mu_t-U\|&=\frac12\sum_x|\mu_t(x)-U(x)|\\&\le\frac{\sqrt n}2\Bigl(\sum_x|\mu_t(x)-U(x)|^2\Bigr)^{1/2}\\&=\frac12\Bigl(\sum_{m\ne0}c_m^{2t}\Bigr)^{1/2}.\end{aligned}
\]
For $1\le m\le n-1$ let $m'=\min(m,n-m)\le\frac n2$; then $c_m=\cos^2\frac{\pi m'}n$, and $\cos u\le e^{-u^2/2}$ for $0\le u\le\frac\pi2$ (the function $\ln\cos u+\frac{u^2}2$ vanishes at $0$
and has derivative $u-\tan u\le0$). Each $m'$ comes from at most two values of $m$, so with $a=\frac{2\pi^2t}{n^2}$,
\[
\begin{aligned}\sum_{m\ne0}c_m^{2t}&\le2\sum_{m'\ge1}e^{-2\pi^2tm'^2/n^2}\eqbr\le2\sum_{m'\ge1}e^{-am'}=\frac{2e^{-a}}{1-e^{-a}}.\end{aligned}
\]
For $t\ge\frac{n^2}4$, $a\ge\frac{\pi^2}2$ and $\frac2{1-e^{-a}}<2.02$, so $d(t)\le\frac12\sqrt{2.02}\,e^{-a/2}\le e^{-\pi^2t/n^2}$.

\emph{Lower bound.} For any $g$ with $|g|\le1$, $|\sum_xg(x)(\mu(x)-\nu(x))|\le\sum_x|\mu(x)-\nu(x)|=2\|\mu-\nu\|$. Take $g(x)=\cos\frac{2\pi x}n$: $\sum_xg(x)U(x)=0$, and $\sum_xg(x)\mu_t(x)$ is the
real part of $\widehat{\mu_t}(1)$, namely $c_1^t$. So $d(t)\ge\frac12\cos^{2t}\frac\pi n$. For $0\le u\le1$, $\cos u\ge1-\frac{u^2}2\ge e^{-u^2}$ (as $1-v\ge e^{-2v}$ for $0\le v\le\frac12$, by
concavity of $v\mapsto1-v-e^{-2v}$, which vanishes at $0$ and is positive at $\frac12$). With $u=\frac\pi n\le1$, $d(t)\ge\frac12e^{-2\pi^2t/n^2}$.

Finally $d(t)>\frac14$ when $t<\frac{\ln2}{2\pi^2}n^2$, and $d(\lceil\frac{n^2}4\rceil)\le e^{-\pi^2/4}<\frac14$.
\end{proof}

\begin{figure}[ht]
\centering
\fig{ch08-cycle}
\caption{The distance $d(t)$ to uniform for the lazy walk on $\Z/n\Z$, for $n=10$, $20$, $40$ and $80$, against $\frac t{n^2}$. The curves almost
coincide: the walk needs a number of steps proportional to $n^2$, the time a random walk takes to travel a distance $n$.}
\label{fig:cycle}
\end{figure}

\section{Shuffling by random transpositions}

A deck of $n$ cards is shuffled as follows: at each step, the left hand picks a uniformly random card $L_t$, the right hand
independently picks a uniformly random card $R_t$, and the two cards are swapped (nothing happens if $L_t=R_t$). This is a
random walk on $\mathfrak S_n$; it is irreducible because transpositions generate $\mathfrak S_n$, aperiodic because it can stay put, and its
stationary law is uniform because each step is a symmetric random walk on a group (the transition matrix is symmetric,
so the uniform law satisfies detailed balance). How many swaps are needed?

To bound mixing times from above we use stopping rules that guarantee perfect mixing. Suppose the chain is driven by
independent random choices $V_1,V_2,\ldots$ (here $V_t=(L_t,R_t)$): $X_t$ is a function of $X_0$ and $V_1,\dots,V_t$. A random time $\tau$ is a
\emph{strong stationary time} if the event $\{\tau=s\}$ is determined by $X_0,V_1,\dots,V_s$ for each $s$, and
for all $s$ and $y$,
\[
P(\tau=s,\ X_s=y)=P(\tau=s)\,\pi(y):
\]
at time $\tau$ the chain is exactly stationary, and independent of $\tau$.

\begin{proposition}\label{prop:sst}
If $\tau$ is a strong stationary time for the chain started at $x$, then $\|P^t(x,\cdot)-\pi\|\le P(\tau>t)$.
\end{proposition}
\begin{proof}
For $s\le t$ and a set $A$, the event $\{\tau=s,\ X_s=z\}$ is determined by $X_0,V_1,\dots,V_s$, and the chain after $s$ is driven by $V_{s+1},\ldots$,
which are independent of it; so $P(\tau=s,\ X_s=z,\ X_t\in A)=P(\tau=s)\pi(z)P^{t-s}(z,A)$. Summing over $z$ and using $\pi P^{t-s}=\pi$,
$P(\tau=s,\ X_t\in A)=P(\tau=s)\pi(A)$. Summing over $s\le t$: $P(X_t\in A)\ge P(\tau\le t,\ X_t\in A)=P(\tau\le t)\pi(A)$, so
$\pi(A)-P(X_t\in A)\le\pi(A)P(\tau>t)\le P(\tau>t)$. The total variation distance is the maximum of $\pi(A)-P^t(x,A)$ over $A$ (take $A=\{\pi>P^t(x,\cdot)\}$).
\end{proof}

\emph{Broder's marking rule.} At the start no card is marked. At step $t$, after the swap, mark $R_t$ if it is unmarked and either
$L_t$ is marked or $L_t=R_t$. Let $M_t$ be the set of marked cards after $t$ steps, and $\tau$ the first time all cards are marked.

\begin{lemma}\label{lem:broder}
For every $t$, conditionally on $M_0,\dots,M_t$, on the set $Q_t$ of positions occupied by marked cards and on the positions of the
unmarked cards, the arrangement of the marked cards in the positions $Q_t$ is uniform among the $|M_t|!$ possibilities.
\end{lemma}
\begin{proof}
By induction on $t$; at $t=0$ there is nothing to arrange. Suppose it holds at time $t$, with $k$ marked cards, and look at
step $t+1$. The choice $(L,R)$ is independent of the past. The sequence $M_0,\dots,M_{t+1}$ is determined by $M_0,\dots,M_t$ and $(L,R)$, since
the marking rule only looks at which cards are chosen. Consider the possible kinds of choice.

\emph{$L$ and $R$ unmarked, $L\ne R$.} The two unmarked cards exchange positions; the marked arrangement is unchanged.

\emph{$L$ and $R$ both marked.} Two marked cards exchange positions: the arrangement is composed with a fixed transposition of
positions, a bijection of the set of arrangements, which preserves the uniform law.

\emph{$L$ unmarked, $R$ marked.} No card is marked; $R$ moves to the position $p'$ of $L$ and $L$ to the position $p$ of $R$. Condition on $L$
and on $p$ (that is, on the event that $R$ is the marked card at position $p$, which has probability $\frac1n$ for each arrangement). The new
positions of the unmarked cards and the new set $Q_{t+1}=(Q_t\setminus\{p\})\cup\{p'\}$ are then fixed, and the new arrangement is the old one
with position $p$ renamed $p'$, a bijection between arrangements: still uniform.

\emph{$R$ unmarked and ($L$ marked or $L=R$): $R$ becomes marked.} Let $q$ be the position of $R$. If $L=R$ nothing moves and $R$ sits at $q$; if $L$ is
the marked card at position $p\in Q_t$, then $R$ goes to $p$ and $L$ to $q$. These $k+1$ cases have probability $\frac1n$ each given $R$. In all of them
$Q_{t+1}=Q_t\cup\{q\}$, the unmarked cards other than $R$ do not move, and the new arrangement of the $k+1$ marked cards is determined by the
old arrangement and by the position of $R$ (one of the $k+1$ positions of $Q_{t+1}$). Conversely the new arrangement determines both:
$R$'s position tells which case occurred, and undoing the swap gives the old arrangement. So the map (old arrangement, case)$\,\mapsto\,$new
arrangement is a bijection from a set of $k!\,(k+1)$ equally likely pairs onto the $(k+1)!$ arrangements, and the new arrangement is
uniform.

In every case, conditionally on the new values of the conditioning variables, the new arrangement is uniform.
\end{proof}

\begin{theorem}[Broder]\label{thm:transpositions}
For random transpositions, $\tau$ is a strong stationary time with $\mathbb E[\tau]=\frac{2n^2}{n+1}H_n\le2nH_n$ and $\operatorname{Var}\tau\le\frac{4\pi^2}3n^2$. Hence for every $c>0$,
\[
d\bigl(2nH_n+cn\bigr)\le\frac{4\pi^2}{3c^2},
\]
and $t_{\mathrm{mix}}\le2n\ln n+O(n)$.
\end{theorem}
\begin{proof}
When all cards are marked, $Q$ is the set of all positions and Lemma~\ref{lem:broder} says that the deck is uniform, conditionally on
$M_0,\dots,M_\tau$, hence on $\tau$: $\tau$ is a strong stationary time ($\{\tau=s\}$ is determined by $M_0,\dots,M_s$, which are functions of the choices).
When $k$ cards are marked, a new card is marked at the next step with probability $p_k=\frac{(n-k)(k+1)}{n^2}$ ($R$ unmarked, and $L$ one of the $k$
marked cards or $L=R$), whatever happened before. So $\tau=\sum_{k=0}^{n-1}G_k$ with independent geometric variables $G_k$ of parameters $p_k$ (the
argument of the coupon collector, Chapter~2). Hence, using $\frac1{(k+1)(n-k)}=\frac1{n+1}\bigl(\frac1{k+1}+\frac1{n-k}\bigr)$,
\[
\mathbb E[\tau]=\sum_{k=0}^{n-1}\frac{n^2}{(k+1)(n-k)}=\frac{n^2}{n+1}\cdot2H_n,
\]
and $\operatorname{Var}\tau=\sum_k\frac{1-p_k}{p_k^2}\le n^4\sum_k\frac1{(k+1)^2(n-k)^2}$. The larger of $k+1$ and $n-k$ is at least $\frac n2$, so $(k+1)(n-k)\ge\frac n2j_k$ with
$j_k=\min(k+1,n-k)$, and each value of $j_k$ occurs for at most two $k$. Hence
\[
\operatorname{Var}\tau\le n^4\cdot\frac4{n^2}\cdot2\sum_{j\ge1}\frac1{j^2}=\frac{4\pi^2}3n^2.
\]
By Chebyshev, $P(\tau>2nH_n+cn)\le P(\tau-\mathbb E\tau>cn)\le\frac{\operatorname{Var}\tau}{c^2n^2}\le\frac{4\pi^2}{3c^2}$, and Proposition~\ref{prop:sst} concludes. As $H_n=\ln n+O(1)$, $t_{\mathrm{mix}}\le2n\ln n+O(n)$.
\end{proof}

\begin{proposition}[lower bound]\label{prop:translower}
For random transpositions with $n\ge2$ and $t\le\frac12n\ln n-cn$ with $c>1$: $d(t)\ge1-6e^{1-2c}$. So $t_{\mathrm{mix}}\ge\frac12n\ln n-2n$.
\end{proposition}
\begin{proof}
Let $F(\sigma)$ be the number of fixed points of a permutation, $A=\{\sigma:F(\sigma)\ge K\}$, and $U_t$ the number of cards never picked by either
hand during the first $t$ steps. Untouched cards are in their original positions, so $F(X_t)\ge U_t$. Under the uniform law,
$\mathbb E[F]=1$ (each card is fixed with probability $\frac1n$), so $\pi(A)\le\frac1K$ by Markov's inequality.

The $2t$ picks are independent uniform cards, so $\mathbb E[U_t]=n(1-\frac1n)^{2t}$, and $\operatorname{Var}U_t\le\mathbb E[U_t]$: writing $U_t=\sum_iI_i$, $\operatorname{Var}I_i\le\mathbb E I_i$ and, for
$i\ne j$, $\operatorname{Cov}(I_i,I_j)=(1-\frac2n)^{2t}-(1-\frac1n)^{4t}\le0$. As $\ln(1-x)\ge-x-x^2$ for $0\le x\le\frac12$ and $\frac{2t}{n^2}\le\frac{\ln n}n\le1$,
\[
\mathbb E[U_t]\ge n\,e^{-2t/n-2t/n^2}\ge n\,e^{-\ln n+2c-1}=e^{2c-1}.
\]
Take $K=\frac12\mathbb E[U_t]$. By Chebyshev, $P(U_t<K)\le\frac{\operatorname{Var}U_t}{K^2}\le\frac4{\mathbb E[U_t]}$. So
\[
\begin{aligned}d(t)&\ge P(X_t\in A)-\pi(A)\eqbr\ge1-\frac4{\mathbb E[U_t]}-\frac2{\mathbb E[U_t]}\ge1-6e^{1-2c}.\end{aligned}
\]
For $c=2$ this is $1-6e^{-3}>0.7$, so $t_{\mathrm{mix}}>\frac12n\ln n-2n$.
\end{proof}

\begin{remark}
The truth is $\frac12n\ln n$: Diaconis and Shahshahani proved in 1981, with the characters of $\mathfrak S_n$ (Tome~2, Chapter~11), that $d(t)$
drops from near $1$ to near $0$ in a window of order $n$ around $\frac12n\ln n$. This sudden drop is the \emph{cutoff phenomenon}. The proof
is given in Tome~9. For $52$ cards, $\frac12n\ln n\approx103$ swaps.
\end{remark}

\section{Riffle shuffles}

The usual way to shuffle is the riffle: cut the deck in two and interleave the halves. The \emph{Gilbert--Shannon--Reeds model}
makes this precise, and the next description generalises it to $a$ packets.

\begin{definition}\label{def:ashuffle}
Number the cards $1,\dots,n$ from top to bottom. An \emph{$a$-shuffle} is performed by choosing a word $w=w_1\cdots w_n\in\{1,\dots,a\}^n$ uniformly
at random (one letter for each position of the shuffled deck). Let $n_j$ be the number of letters equal to $j$. Cut the deck into $a$
packets, the first $n_1$ cards, the next $n_2$, and so on, and fill the positions carrying letter $j$, from top to bottom, with the
cards of the $j$-th packet in their order.
\end{definition}

For $a=2$ this is the GSR riffle shuffle: the cut size $n_1$ is binomial$(n,\frac12)$, and given the cut, all $\binom n{n_1}$ interleavings of the two
packets are equally likely. Equivalently (Exercise~\ref{ex:gsr}), cards drop from the two halves one at a time, from the bottom,
with probability proportional to the current sizes of the halves, which is how a real riffle looks.

A \emph{rising sequence} of an arrangement of the deck is a maximal set of consecutive card numbers $i,i+1,\dots,j$ that appear in
increasing positions. For instance the arrangement $1\,4\,2\,5\,3\,6$ (top to bottom) has the rising sequences $\{1,2,3\}$ and $\{4,5,6\}$.
The rising sequences cut $1,\dots,n$ into blocks, and their number is $r=1+$ the number of $i$ such that card $i+1$ lies above card $i$.

\begin{theorem}[Bayer--Diaconis]\label{thm:bayerdiaconis}
The probability that an $a$-shuffle of the ordered deck produces an arrangement $\sigma$ with $r$ rising sequences is
\[
\frac1{a^n}\binom{a+n-r}n.
\]
\end{theorem}
\begin{proof}
A word $w$ produces $\sigma$ exactly when the cut of $1,\dots,n$ into the packets is compatible with $\sigma$: each packet is a block of consecutive
card numbers, placed in increasing positions, so every $i$ such that card $i+1$ lies above card $i$ in $\sigma$ must be the last card of a
packet. Conversely, given $\sigma$ and packet sizes $n_1,\dots,n_a\ge0$ such that each of these $r-1$ ``descents'' is a packet boundary, there is exactly one
$w$ producing $\sigma$: give letter $j$ to the positions of the cards of packet $j$. Indeed no descent lies inside a packet, so its cards appear in
increasing positions in $\sigma$, and filling the positions of letter $j$ with packet $j$ in order gives back $\sigma$.

So the number of words is the number of ways to choose the $a-1$ boundaries $0\le b_1\le\dots\le b_{a-1}\le n$ (packet $j$ is $b_{j-1}+1,\dots,b_j$, with $b_0=0$, $b_a=n$)
so that each of the $r-1$ descents is among them. Removing one copy of each descent leaves a free multiset of $a-r$ elements of
$\{0,\dots,n\}$, and there are $\binom{(n+1)+(a-r)-1}{a-r}=\binom{a+n-r}n$ of those. (If $r>a$ there is none, and the binomial coefficient is $0$.)
\end{proof}

\begin{proposition}[composition]\label{prop:compose}
An $a$-shuffle followed by an independent $b$-shuffle has the same law as an $ab$-shuffle. So $k$ GSR shuffles are a $2^k$-shuffle.
\end{proposition}
\begin{proof}
Look at the inverse operation. Given the word $w$ and the shuffled deck, the original order is recovered by \emph{sorting stably by the
letters}: put the cards with letter $1$ on top keeping their relative order, then those with letter $2$, and so on. So the inverse of an
$a$-shuffle is: label the cards with independent uniform letters in $\{1,\dots,a\}$, then sort stably by label.

The inverse of ``$a$-shuffle then $b$-shuffle'' is ``inverse $b$-shuffle then inverse $a$-shuffle'': label with uniform $B$-letters and
sort stably, then label with independent uniform $A$-letters and sort stably. Sorting stably by $B$ and then stably by $A$ is sorting
by the pair $(A,B)$ in lexicographic order (this is how radix sort works: after the second sort, cards with different $A$ are in $A$-order, and
cards with the same $A$ keep the order of the first sort, which is $B$-order). The pairs $(A,B)$ are independent and uniform on
$\{1,\dots,a\}\times\{1,\dots,b\}$, which lexicographic order identifies with $\{1,\dots,ab\}$. So the composite inverse is an inverse $ab$-shuffle, and
inverting again, the composite is an $ab$-shuffle.
\end{proof}

Let $A(n,r)$ be the number of permutations of $n$ cards with $r$ rising sequences (the Eulerian numbers). By the theorem and the
proposition, the distance to uniform after $k$ riffle shuffles of $n$ cards is exactly
\[
d_k=\frac12\sum_{r=1}^nA(n,r)\Bigl|\frac1{2^{kn}}\binom{2^k+n-r}n-\frac1{n!}\Bigr|.
\]
The numbers $A(n,r)$ themselves follow from the theorem: the probabilities of all arrangements add up to $1$, so $\sum_rA(n,r)\binom{a+n-r}n=a^n$
for $a=1,\dots,n$, a triangular system (Worpitzky's identity). For $n=52$ the values are:

\begin{center}
\begin{tabular}{c|cccc}
$k$&$1$--$4$&$5$&$6$&$7$\\\hline
$d_k$&$1.000$&$0.924$&$0.614$&$0.334$
\end{tabular}
\qquad
\begin{tabular}{c|ccc}
$k$&$8$&$9$&$10$\\\hline
$d_k$&$0.167$&$0.085$&$0.043$
\end{tabular}
\end{center}

\begin{figure}[ht]
\centering
\fig{ch08-riffle}
\caption{Total variation distance to a uniformly shuffled deck after $k$ riffle shuffles of $52$ cards, computed exactly from the
Bayer--Diaconis formula. It stays near $1$ for four shuffles and then halves with each further shuffle: the cutoff.}
\label{fig:riffle}
\end{figure}

After $k$ shuffles a deck has at most $2^k$ rising sequences, while a uniform permutation of $52$ cards has between
$20$ and $33$ in more than $90\%$ of cases (Exercise~\ref{ex:risingmean}). So with $4$ shuffles, $2^4=16$ rising sequences at most, the deck is visibly far from random,
and a magician can exploit it. Bayer and Diaconis showed that for large $n$ the cutoff is at $\frac32\log_2n$ shuffles; for $n=52$ this
asymptotic formula gives $8.6$, but at this size the drop is spread out, and the table shows it happening between $5$ and $8$.

\begin{beauty}
A card trick, a sorting algorithm and a binomial coefficient: the inverse of a riffle shuffle is one pass of radix sort
on random bits, so $k$ shuffles are one pass on random $k$-bit labels, and the whole law of $k$ shuffles fits in the single
number $\binom{2^k+n-r}n$. Few random processes on a group of size $52!\approx8\times10^{67}$ can be computed exactly; this one can, because of
a hidden structure, the rising sequences.
\end{beauty}

\section{Exercises}

\begin{exercise}[$\star$]\label{ex:lazy}
Show that the lazy chain $\frac12(I+P)$ has the same stationary laws as $P$, is irreducible if $P$ is, and is aperiodic.
\end{exercise}

\begin{exercise}[$\star$]\label{ex:knight}
A knight moves at random on a chessboard, each time choosing uniformly among its legal moves. Show that the expected
number of moves to return to its starting corner is $168$. What is it from a central square?
\end{exercise}

\begin{exercise}[$\star\star$]\label{ex:ehrenfest}
For Ehrenfest's urns with $N$ balls (Example~\ref{ex:examples}), show that the binomial$(N,\frac12)$ law is reversible, and that the expected time
to return to the state ``all balls in the first urn'' is $2^N$. Starting from that state, what does the ergodic theorem say about
the fraction of time spent there? (Boltzmann's answer to Zermelo's objection to the second law of thermodynamics.)
\end{exercise}

\begin{exercise}[$\star\star$]\label{ex:pagerank}
\emph{PageRank.} Let $P$ be any transition matrix on a finite set of $n$ web pages, $0<\alpha<1$, and $G=\alpha P+(1-\alpha)J$, where every row of $J$ is
uniform. Show that $G$ has a unique stationary law $\pi$, that $\|\lambda G^t-\pi\|\le\alpha^t$ for every initial law $\lambda$, and that $\pi=\frac{1-\alpha}n\mathbf 1^{\mathsf T}(I-\alpha P)^{-1}$.
\end{exercise}

\begin{exercise}[$\star\star$]\label{ex:metropolis}
\emph{The Metropolis algorithm.} Let $\pi>0$ be a probability on a finite set and $Q$ a symmetric irreducible transition matrix. Define
$P(x,y)=Q(x,y)\min\bigl(1,\frac{\pi(y)}{\pi(x)}\bigr)$ for $y\ne x$, and $P(x,x)=1-\sum_{y\ne x}P(x,y)$. Show that $P$ is a transition matrix, reversible with respect to
$\pi$, and irreducible. Why is it useful that only the ratios $\frac{\pi(y)}{\pi(x)}$ are needed?
\end{exercise}

\begin{exercise}[$\star\star$]\label{ex:toptorandom}
\emph{Top to random.} A deck is shuffled by taking the top card and inserting it at a uniformly random one of the $n$ possible
places. Let $\tau$ be one step after the time when the card originally at the bottom first reaches the top. Show that $\tau$ is a
strong stationary time, that $\tau$ has the law of the coupon collector's time of Chapter~2, and that $d(n\ln n+cn)\le e^{-c}$.
\end{exercise}

\begin{exercise}[$\star$]\label{ex:gsr}
Show that the GSR shuffle can be performed as follows: cut the deck binomially, then, while both halves are non-empty,
drop the bottom card of a half with probability proportional to its current size onto the shuffled pile.
\end{exercise}

\begin{exercise}[$\star\star$]\label{ex:risingmean}
Show that a uniform random permutation of $n\ge2$ cards has on average $\frac{n+1}2$ rising sequences, with variance $\frac{n+1}{12}$. Deduce from
Chebyshev's inequality that, for $n=52$, more than $90\%$ of all arrangements have between $20$ and $33$ rising sequences.
\end{exercise}

\section{Hints}

\begin{hint}{ex:lazy}
$\pi\frac{I+P}2=\pi$ if and only if $\pi P=\pi$; and $\frac12(I+P)(x,x)\ge\frac12$.
\end{hint}
\begin{hint}{ex:knight}
Corollary~\ref{cor:graph}, and Figure~\ref{fig:knight}.
\end{hint}
\begin{hint}{ex:ehrenfest}
Check detailed balance between $x$ and $x+1$. Kac's formula gives the return time.
\end{hint}
\begin{hint}{ex:pagerank}
$G(x,y)\ge\frac{1-\alpha}n$: apply Lemma~\ref{lem:contraction} with $Q=G$. For the formula, expand $\pi=\pi G$.
\end{hint}
\begin{hint}{ex:metropolis}
Check $\pi(x)P(x,y)=\pi(y)P(y,x)$ when $\pi(y)\le\pi(x)$, and conversely.
\end{hint}
\begin{hint}{ex:toptorandom}
Show by induction that the cards below the original bottom card are in uniformly random order. When $k$ cards are below
it, the top card goes below it with probability $\frac{k+1}n$.
\end{hint}
\begin{hint}{ex:gsr}
Compute the probability of a given interleaving under the dropping rule: the product of the fractions telescopes.
\end{hint}
\begin{hint}{ex:risingmean}
$r=1+\sum_{i<n}D_i$, where $D_i$ is the indicator that card $i+1$ lies above card $i$. Compute $\mathbb E[D_i]$, $\mathbb E[D_iD_{i+1}]$, and note that $D_i$ and $D_j$ are
independent when $|i-j|\ge2$.
\end{hint}

\section{Solutions}

\begin{solution}{ex:lazy}
\idea Laziness only slows the chain down; it cannot change where it spends its time.
\writeup $\pi\cdot\frac12(I+P)=\frac12\pi+\frac12\pi P$, which equals $\pi$ if and only if $\pi P=\pi$. If $P^n(x,y)>0$ then $\bigl(\frac{I+P}2\bigr)^n(x,y)\ge2^{-n}P^n(x,y)>0$ (expand the power and keep
the term $P^n$), so irreducibility is preserved. Finally $\frac{I+P}2(x,x)\ge\frac12>0$, so $1$ belongs to $S_x$ and the period is $1$.
\end{solution}

\begin{solution}{ex:knight}
\idea Do not solve for hitting times: the stationary law of a walk on a graph is proportional to the degrees.
\writeup The knight's graph is connected (a knight can reach every square), so Corollary~\ref{cor:graph} applies. The degrees, listed in Figure~\ref{fig:knight},
add up to $2|\mathcal E|=336$: by rows, $2+3+4+4+4+4+3+2=26$ twice, $3+4+6+6+6+6+4+3=38$ twice, and $4+6+8+8+8+8+6+4=52$ four times, giving
$2\cdot26+2\cdot38+4\cdot52=336$. A corner has degree $2$, so the expected return time is $\frac{336}2=168$. A central square has degree $8$, and the expected
return time is $\frac{336}8=42$.
\end{solution}

\begin{solution}{ex:ehrenfest}
\idea The binomial law is what you get by putting each ball in a random urn independently, and the dynamics moves one
ball at a time without bias.
\writeup With $\pi(x)=\binom Nx2^{-N}$: $\pi(x)P(x,x+1)=\binom Nx\frac{N-x}N2^{-N}$ and $\pi(x+1)P(x+1,x)=\binom N{x+1}\frac{x+1}N2^{-N}$, which are equal since $\binom Nx(N-x)=\binom N{x+1}(x+1)$. So $\pi$ is
reversible, hence stationary (Proposition~\ref{prop:detailed}), and it is the only stationary law as the chain is irreducible. By Kac's formula
the expected return time to $N$ is $\frac1{\pi(N)}=2^N$. The chain has period $2$, but the ergodic theorem does not need aperiodicity: the fraction
of time spent in state $N$ tends to $2^{-N}$, almost surely. For $N=100$, the gas does return to the initial state, but after about $10^{30}$
steps on average; the second law is a statement about overwhelming probabilities.
\end{solution}

\begin{solution}{ex:pagerank}
\idea The damping makes every entry of $G$ at least $\frac{1-\alpha}n$, which is Doeblin's condition with $N=1$.
\writeup $G$ is a transition matrix with $G(x,y)\ge\frac{1-\alpha}n=\delta$, so $|E|\delta=1-\alpha$ and Lemma~\ref{lem:contraction} gives $\|\mu G-\nu G\|\le\alpha\|\mu-\nu\|$. Every entry
is positive, so $G$ is irreducible and aperiodic, and Theorem~\ref{thm:stationary} gives a unique $\pi$. Then $\|\lambda G^t-\pi\|=\|\lambda G^t-\pi G^t\|\le\alpha^t\|\lambda-\pi\|\le\alpha^t$. Since
$\pi J=\frac1n\mathbf 1^{\mathsf T}$ (every row of $J$ is uniform and $\pi$ has mass $1$), $\pi=\pi G$ reads $\pi(I-\alpha P)=\frac{1-\alpha}n\mathbf 1^{\mathsf T}$. The matrix $I-\alpha P$ is invertible because
$\|\alpha\mu P\|_1\le\alpha\|\mu\|_1$ for every row vector $\mu$, so $\mu(I-\alpha P)=0$ forces $\mu=0$. Hence the formula. With $\alpha=0.85$, as in Google's original paper, $50$
iterations of $\lambda\mapsto\lambda G$ give an error below $0.85^{50}<3\times10^{-4}$, whatever the size of the web.
\end{solution}

\begin{solution}{ex:metropolis}
\idea Propose a move with $Q$, then accept it with just the right probability to balance the flows.
\writeup Off the diagonal $0\le P(x,y)\le Q(x,y)$, so $\sum_{y\ne x}P(x,y)\le1$ and $P(x,x)\ge0$: $P$ is a transition matrix. For $x\ne y$, if $\pi(y)\le\pi(x)$ then
$\pi(x)P(x,y)=\pi(x)Q(x,y)\frac{\pi(y)}{\pi(x)}=\pi(y)Q(x,y)$ and $\pi(y)P(y,x)=\pi(y)Q(y,x)\cdot1$, equal as $Q$ is symmetric; the other case is symmetric. So $\pi$ is reversible,
hence stationary. If $Q(x,y)>0$ then $P(x,y)>0$ (as $\pi>0$), so every path of positive $Q$-probability has positive $P$-probability: $P$ is irreducible.
Only ratios are needed, so $\pi$ may be known only up to a constant, as with $\pi(x)\propto e^{-H(x)/T}$ in statistical physics, where the normalising
constant is a sum over astronomically many states. By the ergodic theorem, time averages of the Metropolis chain then compute
averages under $\pi$.
\end{solution}

\begin{solution}{ex:toptorandom}
\idea Every card inserted below the original bottom card lands in a uniformly random place among those already there.
\writeup Call $B$ the card originally at the bottom. We show by induction on $t$ that, conditionally on the number $k$ of cards below $B$ at
time $t$ (and on the history of $k$), the arrangement of those $k$ cards is uniform. A step that does not put the top card below $B$ does not
change them. A step that does inserts the top card in one of the $k+1$ places among them, uniformly, given that it goes below $B$;
inserting a new card in a uniform place of a uniform arrangement of $k$ cards gives a uniform arrangement of $k+1$ cards (each
arrangement arises from exactly one pair). When $B$ reaches the top, the $n-1$ other cards are below it in uniform order, and the
next step inserts $B$ in a uniform place, so at $\tau$ the deck is uniform, independently of $\tau$: a strong stationary time.

When $k$ cards are below $B$, the top card goes below $B$ with probability $\frac{k+1}n$. So the time for $B$ to reach the top is a sum of independent
geometric variables of parameters $\frac1n,\frac2n,\dots,\frac{n-1}n$, and $\tau$ adds one step (a geometric variable of parameter $1$): this is the law of the
coupon collector's time $T_n$ (Chapter~2). For it, $P(T_n>t)\le n(1-\frac1n)^t\le ne^{-t/n}$, since $T_n>t$ means some coupon is still missing after $t$
draws. With $t=n\ln n+cn$ this is $e^{-c}$, and Proposition~\ref{prop:sst} concludes.
\end{solution}

\begin{solution}{ex:gsr}
\idea Both procedures give each pair (cut, interleaving) probability $2^{-n}$.
\writeup Under the dropping rule, given a cut with $n_1$ and $n_2$ cards, the probability of a given interleaving is a product of fractions: when
the halves have $x$ and $y$ cards left, dropping from the first has probability $\frac x{x+y}$. The numerators run through $n_1,n_1-1,\dots,1$ and $n_2,\dots,1$ in
some order, and the denominators through $n,n-1,\dots,1$, so the product is $\frac{n_1!\,n_2!}{n!}=\binom n{n_1}^{-1}$, for every interleaving. Multiplying by the
probability $\binom n{n_1}2^{-n}$ of the cut gives $2^{-n}$, as in Definition~\ref{def:ashuffle} with $a=2$.
\end{solution}

\begin{solution}{ex:risingmean}
\idea A rising sequence ends at each ``descent'' of the inverse permutation, and descents are nearly independent coin flips
of bias $\frac12$.
\writeup Card $i+1$ is above card $i$ with probability $\frac12$ by symmetry, so $\mathbb E[r]=1+\frac{n-1}2=\frac{n+1}2$. The relative order of cards $i,i+1,i+2$ is uniform among
$6$, and only one order (with $i+2$ on top, then $i+1$, then $i$) makes $D_i=D_{i+1}=1$, so $\operatorname{Cov}(D_i,D_{i+1})=\frac16-\frac14=-\frac1{12}$. If $|i-j|\ge2$, $D_i$ and $D_j$ depend on the
relative orders of disjoint pairs of cards, which are independent. So $\operatorname{Var}r=(n-1)\frac14-2(n-2)\frac1{12}=\frac{n+1}{12}$. For $n=52$: mean $26.5$, variance $4.42$, and
Chebyshev gives $P(|r-26.5|\ge7)\le\frac{4.42}{49}<0.1$; as $r$ is an integer, $|r-26.5|<7$ means $20\le r\le33$, which therefore has probability more than $90\%$.
\end{solution}

\section*{Sources and further reading}

\begin{itemize}
\item D.~A.~Levin and Y.~Peres, \emph{Markov Chains and Mixing Times}, 2nd ed., American Mathematical Society, 2017.
Chapters~1--8: stationary laws, coupling, strong stationary times, random transpositions and riffle shuffles, in the
spirit of this chapter.
\item J.~R.~Norris, \emph{Markov Chains}, Cambridge University Press, 1997. A clear account of the classical theory,
including the ergodic theorem by excursions.
\item D.~Bayer and P.~Diaconis, \emph{Trailing the dovetail shuffle to its lair}, Annals of Applied Probability 2 (1992),
294--313. Theorem~\ref{thm:bayerdiaconis} and the seven shuffles.
\item A.~Z.~Broder, \emph{Generating random spanning trees}, FOCS 1989, and P.~Diaconis, \emph{Group Representations in
Probability and Statistics}, IMS, 1988, Chapter~4: strong stationary times and the marking rule.
\item P.~Diaconis and M.~Shahshahani, \emph{Generating a random permutation with random transpositions}, Zeitschrift
für Wahrscheinlichkeitstheorie 57 (1981), 159--179.
\item S.~Brin and L.~Page, \emph{The anatomy of a large-scale hypertextual Web search engine}, Computer Networks and ISDN
Systems 30 (1998), 107--117. PageRank.
\item N.~Metropolis, A.~W.~Rosenbluth, M.~N.~Rosenbluth, A.~H.~Teller and E.~Teller, \emph{Equation of state calculations by
fast computing machines}, Journal of Chemical Physics 21 (1953), 1087--1092.
\item W.~Doeblin, \emph{Exposé de la théorie des chaînes simples constantes de Markov à un nombre fini d'états}, Revue
Mathématique de l'Union Interbalkanique 2 (1938), 77--105. The coupling idea.
\item B.~Hayes, \emph{First links in the Markov chain}, American Scientist 101 (2013), 92--97. Markov's 1913 count of vowels and consonants in \emph{Eugene Onegin}.
\end{itemize}
